\section{Tổng kết và hướng phát triển}
\label{ch:9}

\subsection{Kết quả đạt được của đồ án}

Đồ án tốt nghiệp "Hệ thống Dự báo và Cảnh báo sớm Ùn tắc Giao thông ứng dụng Trí tuệ Nhân tạo" đã được thực hiện và hoàn thành xuất sắc các mục tiêu đề ra ban đầu. Vượt ra khỏi khuôn khổ của một bài toán Học máy (Machine Learning) đơn thuần, nhóm nghiên cứu đã xây dựng thành công một hệ sinh thái công nghệ toàn diện, khép kín (End-to-End Pipeline) bao phủ toàn bộ vòng đời của dữ liệu—từ khâu thu thập, xử lý cho đến việc ra quyết định thông minh và phân phối trực tiếp đến người dùng cuối. Những thành tựu cốt lõi của đồ án trải rộng trên bốn trụ cột chính: Kỹ thuật Dữ liệu, Trí tuệ Nhân tạo, Phát triển Ứng dụng và Vận hành Hệ thống.

\textbf{Về phương diện Kỹ thuật Dữ liệu (Data Engineering):} Hệ thống đã thiết kế và triển khai thành công một Kho dữ liệu (Data Warehouse) chuẩn hóa, đáp ứng trọn vẹn nhu cầu lưu trữ và truy xuất khối lượng lớn dữ liệu giao thông đô thị. Nhóm đã tự động hóa hoàn toàn luồng ETL/ELT, xây dựng một quy trình trích xuất, biến đổi và tải dữ liệu vận hành với độ ổn định cao. Đặc biệt, những khiếm khuyết mang tính đặc thù của dữ liệu cảm biến IoT ngoài thực tế (như mất kết nối, nhiễu tín hiệu) đã được giải quyết triệt để thông qua các kỹ thuật điền khuyết (Forward Fill) và màng lọc nhiễu vật lý, đảm bảo nguồn dữ liệu đầu vào luôn ở trạng thái "sạch" và sẵn sàng cho các mô hình học sâu.

\textbf{Về phương diện Trí tuệ Nhân tạo (Artificial Intelligence):} Đây là điểm sáng học thuật và là đóng góp quan trọng nhất của đồ án. Nhóm đã vượt qua rào cản của bài toán mất cân bằng dữ liệu chuỗi thời gian bằng cách ứng dụng sáng tạo mô hình sinh CTGAN để tổng hợp các kịch bản kẹt xe hiếm gặp, kết hợp cùng màng lọc hậu kiểm vật lý để tạo ra "Tập dữ liệu Vàng". Hơn thế nữa, đồ án đã đề xuất và làm chủ kiến trúc Hybrid Double DQN—một Đặc vụ Học tăng cường (RL Agent) kết hợp mạng Nơ-ron đa nhánh. Nhờ chiến lược Học chuyển giao (Warm-start), Đặc vụ đã triệt tiêu được hội chứng "khởi động lạnh". Sự đột phá lớn nhất nằm ở việc từ bỏ các chỉ số đánh giá tuyến tính thông thường để áp dụng Hàm phần thưởng không đối xứng (Asymmetric Reward Shaping) với "Phạt sinh tử" cho lỗi bỏ lọt thảm họa. Kết quả là mô hình đã vươn lên mốc Độ phủ (Recall) vượt $85\%$ tại các phân lớp kẹt xe nghiêm trọng, hoàn thành triết lý "Thà báo nhầm còn hơn bỏ sót". Cùng với đó, công nghệ XAI (SHAP) đã được tích hợp để giải mã "hộp đen", minh bạch hóa các quyết định cảnh báo dựa trên động lực học dòng chảy vật lý, mang lại độ tin cậy tuyệt đối.

\textbf{Về phương diện Ứng dụng và Định tuyến (Application \& Routing):} Không dừng lại ở các chỉ số trên giấy, năng lực của AI đã được hiện thực hóa thông qua Ứng dụng người dùng (Client-App). Giao diện ứng dụng được thiết kế trực quan, thân thiện (UI/UX), cho phép người dân dễ dàng nắm bắt tình hình giao thông theo thời gian thực. Hệ thống thuật toán định tuyến thông minh đã liên kết trực tiếp với dữ liệu dự báo từ AI, từ đó chủ động cung cấp cho người dùng các lộ trình di chuyển tối ưu nhất nhằm né tránh các điểm nóng kẹt xe ngay cả trước khi chúng kịp hình thành.

\textbf{Về phương diện Triển khai và Vận hành (MLOps \& Deployment):} Toàn bộ các mảnh ghép trên được kết nối thành một Kiến trúc Vận hành chuẩn Công nghiệp (Client-Server). Trái ngược với các nghiên cứu chỉ dừng ở mức độ thử nghiệm (Jupyter Notebook), đồ án đã phân tách rõ ràng các tầng dịch vụ (AI Server, DB Server, App Server). Toàn bộ môi trường, mã nguồn và hệ quy chiếu mô hình đã được container hóa hoàn toàn bằng \textbf{Docker}. Nền tảng CI/CD cơ bản cũng được thiết lập, tạo bệ phóng vững chắc để hệ thống sẵn sàng nhân bản (Scale) và triển khai lên môi trường Đám mây (Cloud) trong thực tế.

\subsection{Những hạn chế còn tồn tại}

Mặc dù đã hoàn thiện một hệ sinh thái End-to-End mạnh mẽ, "Hệ thống Dự báo và Cảnh báo sớm Ùn tắc Giao thông" vẫn không tránh khỏi những điểm giới hạn nhất định. Nguyên nhân chủ yếu xuất phát từ rào cản tài nguyên phần cứng, giới hạn thời gian nghiên cứu và bản chất cực kỳ phức tạp của mạng lưới giao thông đô thị.

\textbf{Thứ nhất, giới hạn về Nhận thức Không gian và sự Lan truyền ùn tắc:} Dù mạng Q-Network hiện tại sở hữu khả năng nắm bắt động lực học chuỗi thời gian xuất sắc nhờ LSTM, nó vẫn xử lý các phân đoạn đường bộ một cách tương đối độc lập thông qua Embedding ID. Hệ thống chưa mô hình hóa được cấu trúc mạng lưới giao thông thành một đồ thị toán học (Graph), khiến việc nắm bắt "hiệu ứng dội ngược" (Spillback) hoặc sự lan truyền kẹt xe giữa các nút giao liền kề vẫn mang tính gián tiếp. Bên cạnh đó, cơ chế dự phòng (Global Spatial Fallback) khi mất tín hiệu cảm biến mới chỉ dừng ở mức xấp xỉ thống kê không gian, chưa phản ánh trọn vẹn luồng di chuyển vật lý của dòng xe.

\textbf{Thứ hai, khoảng cách giữa Không gian học và Thực tiễn (Sim-to-Real Gap):} Tập dữ liệu huấn luyện dù đã được làm giàu nhưng vẫn duy trì độ "sạch" nhất định. Khi phải đối mặt trực tiếp với môi trường thực tế—nơi đầy rẫy các "sự kiện thiên nga đen" (Black Swan) như tai nạn bất ngờ, mất kết nối thiết bị hay sai số GPS nhảy vọt—Đặc vụ (Agent) có nguy cơ bị "sốc" do chưa từng quan sát thấy mức nhiễu loạn cực đoan này trong không gian Offline RL. Hơn nữa, việc mạng Nơ-ron bị đóng băng trọng số sau khi triển khai khiến hệ thống thiếu đi khả năng tự học liên tục (Continuous Learning) từ những sai lầm mới phát sinh.

\textbf{Thứ ba, rào cản hiệu năng ở quy mô siêu đô thị:} Toàn bộ pipeline suy luận hiện tại đang được thực thi trên môi trường Python/PyTorch tiêu chuẩn. Tuy hoạt động mượt mà ở quy mô thử nghiệm, nhưng nếu mở rộng ra hàng vạn đoạn đường cùng lúc, sự hạn chế về đa luồng của Python (GIL) và kiến trúc chưa được tối ưu hóa cho phần cứng sẽ tạo ra nút thắt cổ chai về độ trễ (Latency), thách thức trực tiếp tiêu chuẩn khắt khe của hệ thống cảnh báo thời gian thực.

\textbf{Cuối cùng, gánh nặng chi phí hạ tầng (Data Streaming Infrastructure):} Để duy trì tính liên tục của dữ liệu, việc Polling liên tục từ các API hoặc duy trì kết nối luồng dữ liệu lớn đòi hỏi chi phí băng thông và máy chủ đám mây khổng lồ. Để cân đối ngân sách, hệ thống đôi khi phải chấp nhận độ trễ thông qua xử lý lô nhỏ (batch-processing 10-15 phút). Đồng thời, thuật toán tái định tuyến trên App vẫn còn độ trễ cập nhật nhất định nếu một thảm họa kẹt xe xảy ra ngay khi người dùng đang ở giữa hành trình.

\subsection{Hướng phát triển trong tương lai}

Nhận thức rõ những thách thức trên, nhóm nghiên cứu đã phác thảo một lộ trình nâng cấp toàn diện, định hướng hệ thống tiến xa hơn trên con đường trở thành một nền tảng quản lý giao thông thông minh cốt lõi của đô thị.

\textbf{Về đột phá Kiến trúc Trí tuệ Nhân tạo:} Đích đến tiếp theo sẽ là sự tích hợp của Mạng Nơ-ron Đồ thị (Graph Neural Networks - GNN). Bằng cách chuyển đổi mô hình không gian từ chuỗi độc lập sang cấu trúc đồ thị động $G=(V, E)$ thông qua các kiến trúc như STGCN hay GAT, Đặc vụ sẽ có được "tầm nhìn chim bay" bao quát toàn bộ mạng lưới, đo lường chính xác hiệu ứng lan truyền của từng vụ kẹt xe. Đồng thời, hệ thống sẽ được chuyển dịch sang mô hình Học tăng cường Trực tuyến (Online RL). Cơ chế vòng kín này cho phép Đặc vụ liên tục đối chiếu dự báo với thực tế, sử dụng sai số để tinh chỉnh trọng số liên tục theo thời gian thực, tạo ra một "hệ miễn dịch học tập" thích ứng hoàn hảo với sự thay đổi của cơ sở hạ tầng.

\textbf{Về tối ưu hóa Hiệu năng và Vận hành (MLOps):} Bài toán "nút thắt cổ chai" độ trễ sẽ được giải quyết triệt để thông qua việc xuất mạng Q-Network sang định dạng ONNX và biên dịch bằng NVIDIA TensorRT. Áp dụng kỹ thuật lượng tử hóa (Quantization) từ FP32 xuống FP16/INT8 sẽ giúp gia tốc suy luận lên hàng chục lần, đáp ứng mượt mà SLA thời gian thực. Hơn thế nữa, kiến trúc Điện toán Biên (Edge Computing) sẽ được triển khai. Việc đẩy các mô hình AI hạng nhẹ xuống trực tiếp các thiết bị biên tại ngã tư, kết hợp nền tảng Data Streaming như Apache Kafka, sẽ san sẻ gánh nặng cho hệ thống Cloud trung tâm, tạo tiền đề để mở rộng quy mô toàn thành phố một cách bền vững.

\textbf{Về Nâng cấp Ứng dụng và Định tuyến:} Ứng dụng người dùng sẽ được nâng cấp khả năng "Định tuyến động trong hành trình" (Dynamic In-trip Re-routing). Thông qua kết nối WebSockets thời gian thực, ứng dụng sẽ liên tục "lắng nghe" luồng cảnh báo từ AI. Bất cứ khi nào phát hiện rủi ro phía trước, thuật toán sẽ ngay lập tức tính toán và gợi ý chuyển hướng ngay cả khi phương tiện đang lăn bánh. Đặc biệt, hệ thống sẽ khai phá sức mạnh của Dữ liệu Đám đông (Crowdsourcing). Bằng cách cho phép người dùng chủ động báo cáo các sự cố trên đường, nguồn tin cộng đồng kết hợp cùng dữ liệu IoT sẽ tạo thành một tấm lưới cảm biến khổng lồ, đưa độ nhạy và tính chính xác của Hệ thống Cảnh báo Sớm lên một tầm cao mới.
